
\documentclass{abgabe}

\usepackage{textcomp}  % Required for encoding \textbigcircle
\usepackage{scalerel}  % Required for emoji \scalerel

\def\☀{\scalerel*{\includegraphics{figures/1f506.eps}}{\textrm{\textbigcircle}}}
\def\🦵{\scalerel*{\includegraphics{figures/1f9b5.eps}}{\textrm{\textbigcircle}}}

\begin{document}

\begin{questions}
    \question
    Von 20 Teilnehmern einer Bergwanderung geben 8 Personen an Knieschmerzen zu haben.
    6 Teilnehmer leiden unter Sonnenbrand.
    8 Teilnehmer bleiben unversehrt.

    \begin{center}
        \begin{venndiagram2sets}[labelA={\🦵},labelB={\☀},labelOnlyA={8?},labelOnlyB={6?},labelNotAB={8}]
            %\fillNotA
        \end{venndiagram2sets}
    \end{center}

    Wie groß ist die Wahrscheinlichkeit, dass ein zufällig ausgewählter Teilnehmer
    \begin{parts}
        \part
        Sonnenbrand oder Knieschmerzen hat?
        \begin{solution}
            Es gilt:
            \[
                P(\☀ \cup \🦵) = 1 - P(\conj{\☀ \cup \🦵}) = 1 - \frac{8}{20} = \frac{12}{20} = 0.6
            \]
            \qed
        \end{solution}

        \part
        Knieschmerzen und Sonnenbrand aufweist?
        \begin{solution}
            Es gilt:
            \[
                P(\☀ \cap \🦵) = P(\☀) + P(\🦵) - P(\☀ \cup \🦵) = \frac{6}{20} + \frac{8}{20} - \frac{12}{20} = \frac{2}{20} = 0.1
            \]

            \qed
        \end{solution}

        \part
        der Knieschmerzen hat, keinen Sonnenbrand hat?
        \begin{solution}
            Es gilt:
            \[
                P(\conj{\☀} \cap \🦵) = P(\🦵) - P(\☀ \cap \🦵) = \frac{8}{20} - \frac{2}{20} = \frac{6}{20} = 0.3
            \]
            \qed
        \end{solution}

        \part
        Sonnenbrand, aber keine Knieschmerzen hat?
        \begin{solution}
            Es gilt:
            \[
                P(\☀ \cap \conj{\🦵}) = P(\☀) - P(\☀ \cap \🦵) = \frac{6}{20} - \frac{2}{20} = \frac{4}{20} = 0.2
            \]
            \qed
        \end{solution}
    \end{parts}
\end{questions}
\end{document}